\documentclass[12pt]{article}

\usepackage[margin=1in]{geometry}
\usepackage{amsmath, amssymb}
\usepackage{fancyhdr}

\pagestyle{fancy}
\fancyhf{}
\cfoot{\thepage}
\renewcommand{\headrulewidth}{0pt}

\begin{document}

\begin{center}

    {\Large \textbf{02-680 Essential Mathematics and Statistics}}
    
    \vspace{0.4em}
    
    {\large \textbf{Homework 1}}
    
    \vspace{0.7em}
    

    
    \vspace{0.2em}
    
    Andrew ID: 
    
    \vspace{0.2em}
    
    September 12, 2026

\end{center}

\vspace{1em}


\vspace{1em}

\vspace{1em}

%%%%%%%%%%%%%%%%%%%%%%%%%%%%%%%%%%%%%%%%%%%
%%%%%%%%%%%%%%%%%%%%%%%%%%%%%%%%%%%%%%%%%%%
%%%%%%%%%%%%%%%%%%%%%%%%%%%%%%%%%%%%%%%%%%%
\textbf{Problem 1. Logics and Sets}

\begin{enumerate}
     \item[(a)]
     \begin{enumerate}
        \item[i.]
        
        The median operator is true when at least two of $p,q,r$ are true. Thus, an equivalent compound proposition is:
        \[
        (\neg p\land q \land r)\lor(p\land \neg q \land r)\lor(p\land q \land \neg r)\lor(p\land q \land r)
        \]
        
        Compare the expression with the given truth table:
        \[
        \begin{array}{ccc|c|c}
        p&q&r&p\nabla q\nabla r&(\neg p\land q \land r)\lor(p\land \neg q \land r)\lor(p\land q \land \neg r)\lor(p\land q \land r)\\
        \hline
        T&T&T&T&T\\
        T&T&F&T&T\\
        T&F&T&T&T\\
        T&F&F&F&F\\
        F&T&T&T&T\\
        F&T&F&F&F\\
        F&F&T&F&F\\
        F&F&F&F&F
        \end{array}
        \]
        The expression macthes every row.
        
        \item[ii.]
        
        An element $x$ belongs to $M(A,B,C)$ exactly when
        $x$ belongs to at least two of the sets $A,B,C$. Therefore,
        \[
        M(A,B,C)
        =
        (A\cap B \cap \overline{C})\cup(A\cap \overline{B} \cap C)\cup(\overline{A}\cap B \cap C)\cup(A\cap B \cap C).
        \]
        
        Correspondence: $\neg p\land q \land r$ corresponds to $\overline{A}\cap B \cap C$, $p\land \neg q \land r$ corresponds to $A\cap \overline{B} \cap C$...... The $\lor$ corresponds to $\cup$.
        
        \item[iii.]

        We want to prove that
        \[
        \begin{aligned}
        &(A\cap B\cap\overline{C})
        \cup(A\cap\overline{B}\cap C)
        \cup(\overline{A}\cap B\cap C)
        \cup(A\cap B\cap C)\\
        &\qquad=
        (A\cap B)\cup(A\cap C)\cup(B\cap C).
        \end{aligned}
        \]
        
        We simplify the left-hand side (LHS):
        \begin{align*}
        \text{LHS}
        &= (A\cap B\cap\overline{C})
           \cup(A\cap B\cap C)
           \cup(A\cap\overline{B}\cap C)
           \cup(\overline{A}\cap B\cap C)\\
        &= \big((A\cap B)\cap(\overline{C}\cup C)\big)
           \cup\big((A\cap C)\cap(\overline{B}\cup B)\big)
           \cup\big((B\cap C)\cap(\overline{A}\cup A)\big).
        \end{align*}
        
        For any set $X$, $X\cup\overline{X}=U$, and $X\cap U=X$. Therefore,
        \begin{align*}
        \text{LHS}
        &= \big((A\cap B)\cap U\big)
           \cup\big((A\cap C)\cap U\big)
           \cup\big((B\cap C)\cap U\big)\\
        &= (A\cap B)\cup(A\cap C)\cup(B\cap C)\\
        &= \text{RHS}.
        \end{align*}
        
        Thus we can conclude that,
        \[
        M(A,B,C)
        =
        (A\cap B)\cup(A\cap C)\cup(B\cap C).
        \]
    \end{enumerate}

%%%%%%%%%%%%%%%%%%%%%%%%%%%%%%%%%%%%%%%%%%%
%%%%%%%%%%%%%%%%%%%%%%%%%%%%%%%%%%%%%%%%%%%
    \item[(b)]
    \begin{enumerate}
        \item [i.]
        A gene $g\in R_{\text{active}}$ is defined by $\exists h\in G: R(g,h)$.
        $G\setminus R_{\text{active}}$ is the negation:
        \[
        \neg\big(\exists h\in G: R(g,h)\big)
        \]
        which equals to:
        \[
        \forall h\in G: \neg R(g,h).
        \]
        For every gene $h\in G$, gene $g$ does not repress $h$.
        
        \item[ii.]
        Let $G=\{a,b,c\}$.
        Repression relations: only $R(a,b)$ holds.\\
        $a$ represses $b$, so $a\in R_{\text{active}}$.\\
        $b$ represses no genes, so $b\in G\setminus R_{\text{active}}$.\\
        $c$ represses no genes, so $c\in G\setminus R_{\text{active}}$.\\
        We have $R_{\text{active}}=\{a\}$, so $R_{\text{active}}\neq G$.
        Genes in $G\setminus R_{\text{active}}$: $b,c$.
        
        \item[iii.]
        Yes, it is possible.

        Example: $G=\{x\}$, a network with exactly one gene, and no repression relations.\\
        Thus $R_{\text{active}}=\emptyset$.
    \end{enumerate}
    
    
\end{enumerate}

\vspace{1em}


%%%%%%%%%%%%%%%%%%%%%%%%%%%%%%%%%%%%%%%%%%%
%%%%%%%%%%%%%%%%%%%%%%%%%%%%%%%%%%%%%%%%%%%
%%%%%%%%%%%%%%%%%%%%%%%%%%%%%%%%%%%%%%%%%%%
\textbf{Problem 2. Sets, Functions, Tuples, and Asymptotics}

\begin{enumerate}
    \item [(a)]
    \begin{enumerate}
        \item [i.]
        Choose $n_1$ freely: $4$ choices. Once $n_1$ is fixed, $n_2$ may be any of the remaining
        $3$ bases.
        \[
        |S|=4\cdot 3=12.
        \]
        \item[ii.] 
        Yes. For example,
        \[
        f(\langle A,C\rangle)=C,\qquad f(\langle G,C\rangle)=C,
        \]
        For every $n_2\in\Sigma$, choose any $n_1\in\Sigma$ with $n_1\neq n_2$; then
        $\langle n_1,n_2\rangle\in S$ and $f$ maps it to $n_2$. Hence every element of
        $\Sigma$ is attained, so
        \[
        \operatorname{range}(f)=\Sigma.
        \]
        Thus the range equals the codomain.
    \end{enumerate}
%%%%%%%%%%%%%%%%%%%%%%%%%%%%%%%%%%%%%%%%%%%
%%%%%%%%%%%%%%%%%%%%%%%%%%%%%%%%%%%%%%%%%%%
    \item[(b)]
    \begin{enumerate}
        \item [i.]
        Build the tuple left to right: $n_1$ has $4$ choices, and for each $i\ge 2$, $n_i$
        only needs to differ from $n_{i-1}$, giving 3 choices.
        Therefore
        \[
        {c(k)=4\cdot 3^{\,k-1}\qquad (k\ge 1).}
        \]
        
        \item[ii.]
        \[
        c(k)=4\cdot 3^{k-1}=\frac{4}{3}\cdot 3^k,
        \]
        
        For all $k\ge 1$:
        Taking $c=10$ gives $c(k)\le c\cdot 3^k$, so $c(k)\in O(3^k)$.
        Taking $c'=1$ gives $c(k)\ge c'\cdot 3^k$, so $c(k)\in\Omega(3^k)$.
        
        Thus,
        \[
        {c(k)\in\Theta(3^k).}
        \]
    \end{enumerate}
\end{enumerate}

%%%%%%%%%%%%%%%%%%%%%%%%%%%%%%%%%%%%%%%%%%%
%%%%%%%%%%%%%%%%%%%%%%%%%%%%%%%%%%%%%%%%%%%
%%%%%%%%%%%%%%%%%%%%%%%%%%%%%%%%%%%%%%%%%%%
\textbf{Problem 3. Sets, Functions, and Tuples}

\begin{enumerate}
    \item [(a)]
    \begin{enumerate}
        \item [i.]
        Given $D_X = \{g_1,g_2,g_3,g_5,g_7\}$ and $D_Y = \{g_2,g_3,g_4,g_6,g_7,g_8\}$:
        \[
        D_X\cap D_Y = \{g_2,g_3,g_7\}
        \]
        \[
        D_X\cup D_Y = \{g_1,g_2,g_3,g_4,g_5,g_6,g_7,g_8\}
        \]
        
        \item[ii.]
        \[
        \begin{array}{cc|c}
        \chi_X(g) & \chi_Y(g) & \chi_X(g)\oplus\chi_Y(g) \\
        \hline
        0 & 0 & 0 \\
        0 & 1 & 1 \\
        1 & 0 & 1 \\
        1 & 1 & 0
        \end{array}
        \]
        Condition-specific genes: $\{g_1,g_4,g_5,g_6,g_8\}$.

        Set-builder expression:
        \[
        (D_X \setminus D_Y)\cup(D_Y\setminus D_X)
        \]
        Verification:
        $D_X\setminus D_Y=\{g_1,g_5\}$, $D_Y\setminus D_X=\{g_4,g_6,g_8\}$.
        Their union $\{g_1,g_4,g_5,g_6,g_8\}$ gives exactly the same list.
    \end{enumerate}
\end{enumerate}





%%%%%%%%%%%%%%%%%%%%%%%%%%%%%%%%%%%%%%%%%%%
%%%%%%%%%%%%%%%%%%%%%%%%%%%%%%%%%%%%%%%%%%%
%%%%%%%%%%%%%%%%%%%%%%%%%%%%%%%%%%%%%%%%%%%
\textbf{Problem 4. Logic, Functions, and Asymptotics}

\begin{enumerate}
    \item [(a)]
    \begin{enumerate}
        \item [i.]
        \[
        f(x) \in O(g(x)) \iff \exists x_0 \ge 0,\ \exists c > 0:\ \forall x \ge x_0,\ f(x) \le c\cdot g(x).
        \]
        
        The negation of the right-hand side:
        \begin{align*}
        \neg\Big(\exists x_0 \ge 0,\ \exists c > 0:\ \forall x \ge x_0,\ f(x) \le c g(x)\Big)
        &= \forall x_0 \ge 0,\ \neg\Big(\exists c > 0:\ \forall x \ge x_0,\ f(x) \le c g(x)\Big) \\
        &= \forall x_0 \ge 0,\ \forall c > 0:\ \neg\Big(\forall x \ge x_0,\ f(x) \le c g(x)\Big) \\
        &= \forall x_0 \ge 0,\ \forall c > 0:\ \exists x \ge x_0,\ \neg\big(f(x) \le c g(x)\big) \\
        &= \forall x_0 \ge 0,\ \forall c > 0:\ \exists x \ge x_0,\ f(x) > c g(x).
        \end{align*}
        So $f(x)\notin O(g(x))$ means:
        \[
        {\forall x_0 \ge 0,\ \forall c > 0:\ \exists x \ge x_0,\ f(x) > c\cdot g(x)}.
        \]

        
        \item[ii.]

        Take any arbitrary $x_0\ge 0$ and any arbitrary constant $c>0$.
        To prove $x^2 \notin O(x)$, we need to pick an $x \ge x_0$ such that $x^2 > c x$.
        
        Solve $x^2 > c x$. Since $x>0$, divide both sides by $x$:
        \[
        x > c.
        \]
        Choose $x = \max(x_0,\ c+1)$. It satifies that $x^2 > c x$.


        By the negation from i., $x^2 \notin O(x)$.
    \end{enumerate}
%%%%%%%%%%%%%%%%%%%%%%%%%%%%%%%%%%%%%%%%%%%
%%%%%%%%%%%%%%%%%%%%%%%%%%%%%%%%%%%%%%%%%%%

    \item [(b)]
    \begin{enumerate}
        \item [i.]
        There exists constants $n_0\ge1$, $k=10$ such that
        \[
        \forall n\ge n_0,\  n \le k\cdot n.
        \]
        Thus
        \[
        \quad n\in O(n).
        \]

        We know:
        \[
        \log n \in O(n).
        \]
        Apply the product-rule property:
        \[
        (\log n)\cdot n \in O(n\cdot n) = O(n^2).
        \]
        Therefore $r(n)=n\log n \in O(n^2)$.
        
        \item[ii.]        
        From i.: $n\log n \in O(n^2)$.
        This means that $r(n)$ grows slower than or equal to $p(n)$.

        
        $p(n)=n^2$ is a polynomial, so $n^2\in O(2^n)$.
        Thus $p(n)$ grows slower than or equal to $q(n)$.
        
        Order from slowest-growing to fastest-growing:
        \[
        {r(n)=n\log n,\quad p(n)=n^2,\quad q(n)=2^n}.
        \]
        The O relationships between each consecutive pair:
        \[
        n\log n \in O(n^2),  \quad  n^2\in O(2^n).
        \]
    \end{enumerate}
\end{enumerate}

\end{document}
